% --- Archon physics formalization source begin ---
% archon:physics
% archon:covers PhyXMiniProblems/problem_phyx_mini_0931.lean
% archon:source-report reports/phyx_mini/problem_phyx_mini_0931.source.json

\chapter{Physics problem phyx\_mini\_0931}
\label{ch:PhyXMiniProblems_problem_phyx_mini_0931}

\paragraph{Problem source.}
\#\# Physical scenario

The square loop shown in figure moves into a \textbackslash{}( 0.80 \textbackslash{}, \textbackslash{}text\{T\} \textbackslash{}) magnetic field at a constant speed of \textbackslash{}( 10 \textbackslash{}, \textbackslash{}text\{m/s\} \textbackslash{}). The loop has a resistance of \textbackslash{}( 0.10 \textbackslash{}, \textbackslash{}Omega \textbackslash{}), and it enters the field at \textbackslash{}( t = 0 \textbackslash{}, \textbackslash{}text\{s\} \textbackslash{}).

\#\# Question

What is the maximum current?

\#\# Answer choices

- A: A: \textbackslash{}( 13 \textbackslash{}

- B: B: \textbackslash{}( 11 \textbackslash{}

- C: C: \textbackslash{}( 10 \textbackslash{}

- D: D: \textbackslash{}( 9 \textbackslash{}

\#\# Figure caption (auxiliary; use the image as primary evidence)

The image features a square loop tilted at an angle, resembling a diamond shape. The loop is outlined in orange. Each side of the loop is labeled as "10 cm". 

To the right of the loop, there is a dashed vertical line marking a boundary. To the right of this boundary, there are numerous blue "X" symbols arranged in a grid pattern, indicating a magnetic field. The magnetic field is directed into the plane of the image.

Next to the field region, it is labeled with "B = 0.80 T," specifying the magnetic field strength.

A green arrow points to the right, originating from the center of the loop, with the label "10 m/s," indicating the velocity of the loop moving towards the magnetic field region.

\#\# Dataset metadata

Electromagnetic Induction; Physical Model Grounding Reasoning, Multi-Formula Reasoning

\paragraph{Recorded answer/context.}
B: B: \textbackslash{}( 11 \textbackslash{}

\paragraph{Figure/image path.}
/root/proposal\_for\_physic/science-mango/phyx\_mini\_run/phyx\_data/test\_image/931.png

\paragraph{Formalization target.}
create a compiling Lean file with sorry bodies at `PhyXMiniProblems/problem\_phyx\_mini\_0931.lean`. The Lean declarations must preserve the physical quantities, dimensions or dimensional roles, figure labels, governing-law hypotheses, and final relation expressed by this problem.
Use Mathlib/Physlib names found through LeanExplore where available. If a physics API is missing, introduce faithful local abstractions rather than scalar placeholder aliases.

\begin{theorem}[Physics formalization target]
\label{thm:physics:phyx_mini_0931:target}
\lean{PhyXMiniProblems.ProblemPhyXMini0931.problem_phyx_mini_0931}
\uses{def:physics:phyx-mini-0931:phyxminiproblems-problemphyxmini0931-movingsquareloopsetup, def:physics:phyx-mini-0931:phyxminiproblems-problemphyxmini0931-matchesmovingsquareloopscenario, def:physics:phyx-mini-0931:phyxminiproblems-problemphyxmini0931-matchessuppliedmovingsquareloopfigure, def:physics:phyx-mini-0931:phyxminiproblems-problemphyxmini0931-matchesproblemmeasurements, def:physics:phyx-mini-0931:phyxminiproblems-problemphyxmini0931-hasphysicalmovingloopparameters, def:physics:phyx-mini-0931:phyxminiproblems-problemphyxmini0931-modelsuniformintopagehalfplanefield, def:physics:phyx-mini-0931:phyxminiproblems-problemphyxmini0931-satisfiesconstanttranslationkinematics, def:physics:phyx-mini-0931:phyxminiproblems-problemphyxmini0931-satisfiesdiamondboundarysweepgeometry, def:physics:phyx-mini-0931:phyxminiproblems-problemphyxmini0931-satisfiesuniformoverlapfluxlaw, def:physics:phyx-mini-0931:phyxminiproblems-problemphyxmini0931-satisfiesfaradayinductionlaw, def:physics:phyx-mini-0931:phyxminiproblems-problemphyxmini0931-satisfiesohmslaw, def:physics:phyx-mini-0931:phyxminiproblems-problemphyxmini0931-ismaximumcurrentinamperes, def:physics:phyx-mini-0931:phyxminiproblems-problemphyxmini0931-recordeddatasetanswer, def:physics:phyx-mini-0931:phyxminiproblems-problemphyxmini0931-isclosestdisplayedcurrent}
The assigned autoformalize agent should translate this physics problem into Lean declarations in the covered file, with theorem and lemma proof bodies written as `by sorry`.
\end{theorem}
\begin{proof}
This is an autoformalization task, not a proof task. Produce statements that can later be proved by the physics prover without weakening the physical model.
\end{proof}

% --- Archon declaration topology begin ---
\section{Declaration topology}
\begin{definition}[area Dimension]
  \label{def:physics:phyx-mini-0931:phyxminiproblems-problemphyxmini0931-areadimension}
  \lean{PhyXMiniProblems.ProblemPhyXMini0931.areaDimension}
  Physical dimension L² of area
\end{definition}
\begin{proof}
  This is the declared model definition of area Dimension.
\end{proof}
\begin{definition}[area Rate Dimension]
  \label{def:physics:phyx-mini-0931:phyxminiproblems-problemphyxmini0931-arearatedimension}
  \lean{PhyXMiniProblems.ProblemPhyXMini0931.areaRateDimension}
  \uses{def:physics:phyx-mini-0931:phyxminiproblems-problemphyxmini0931-areadimension}
  Physical dimension L² T⁻¹ of an area-change rate
\end{definition}
\begin{proof}
  This is the declared model definition of area Rate Dimension.
\end{proof}
\begin{definition}[magnetic Flux Density Dimension]
  \label{def:physics:phyx-mini-0931:phyxminiproblems-problemphyxmini0931-magneticfluxdensitydimension}
  \lean{PhyXMiniProblems.ProblemPhyXMini0931.magneticFluxDensityDimension}
  Physical dimension M T⁻¹ C⁻¹ of magnetic flux density
\end{definition}
\begin{proof}
  This is the declared model definition of magnetic Flux Density Dimension.
\end{proof}
\begin{definition}[magnetic Flux Dimension]
  \label{def:physics:phyx-mini-0931:phyxminiproblems-problemphyxmini0931-magneticfluxdimension}
  \lean{PhyXMiniProblems.ProblemPhyXMini0931.magneticFluxDimension}
  \uses{def:physics:phyx-mini-0931:phyxminiproblems-problemphyxmini0931-areadimension, def:physics:phyx-mini-0931:phyxminiproblems-problemphyxmini0931-magneticfluxdensitydimension}
  Physical dimension M L² T⁻¹ C⁻¹ of magnetic flux
\end{definition}
\begin{proof}
  This is the declared model definition of magnetic Flux Dimension.
\end{proof}
\begin{definition}[electromotive Force Dimension]
  \label{def:physics:phyx-mini-0931:phyxminiproblems-problemphyxmini0931-electromotiveforcedimension}
  \lean{PhyXMiniProblems.ProblemPhyXMini0931.electromotiveForceDimension}
  \uses{def:physics:phyx-mini-0931:phyxminiproblems-problemphyxmini0931-magneticfluxdimension}
  Physical dimension M L² T⁻² C⁻¹ of electromotive force
\end{definition}
\begin{proof}
  This is the declared model definition of electromotive Force Dimension.
\end{proof}
\begin{definition}[electric Current Dimension]
  \label{def:physics:phyx-mini-0931:phyxminiproblems-problemphyxmini0931-electriccurrentdimension}
  \lean{PhyXMiniProblems.ProblemPhyXMini0931.electricCurrentDimension}
  Physical dimension C T⁻¹ of electric current
\end{definition}
\begin{proof}
  This is the declared model definition of electric Current Dimension.
\end{proof}
\begin{definition}[electrical Resistance Dimension]
  \label{def:physics:phyx-mini-0931:phyxminiproblems-problemphyxmini0931-electricalresistancedimension}
  \lean{PhyXMiniProblems.ProblemPhyXMini0931.electricalResistanceDimension}
  Physical dimension M L² T⁻¹ C⁻² of electrical resistance
\end{definition}
\begin{proof}
  This is the declared model definition of electrical Resistance Dimension.
\end{proof}
\begin{definition}[Length Quantity]
  \label{def:physics:phyx-mini-0931:phyxminiproblems-problemphyxmini0931-lengthquantity}
  \lean{PhyXMiniProblems.ProblemPhyXMini0931.LengthQuantity}
  A nonnegative, unit-independent physical length
\end{definition}
\begin{proof}
  This is the declared model definition of Length Quantity.
\end{proof}
\begin{definition}[Time Quantity]
  \label{def:physics:phyx-mini-0931:phyxminiproblems-problemphyxmini0931-timequantity}
  \lean{PhyXMiniProblems.ProblemPhyXMini0931.TimeQuantity}
  A nonnegative, unit-independent clock time
\end{definition}
\begin{proof}
  This is the declared model definition of Time Quantity.
\end{proof}
\begin{definition}[Speed Quantity]
  \label{def:physics:phyx-mini-0931:phyxminiproblems-problemphyxmini0931-speedquantity}
  \lean{PhyXMiniProblems.ProblemPhyXMini0931.SpeedQuantity}
  A nonnegative, unit-independent speed
\end{definition}
\begin{proof}
  This is the declared model definition of Speed Quantity.
\end{proof}
\begin{definition}[Area Quantity]
  \label{def:physics:phyx-mini-0931:phyxminiproblems-problemphyxmini0931-areaquantity}
  \lean{PhyXMiniProblems.ProblemPhyXMini0931.AreaQuantity}
  \uses{def:physics:phyx-mini-0931:phyxminiproblems-problemphyxmini0931-areadimension}
  A nonnegative, unit-independent area
\end{definition}
\begin{proof}
  This is the declared model definition of Area Quantity.
\end{proof}
\begin{definition}[Area Rate Magnitude Quantity]
  \label{def:physics:phyx-mini-0931:phyxminiproblems-problemphyxmini0931-arearatemagnitudequantity}
  \lean{PhyXMiniProblems.ProblemPhyXMini0931.AreaRateMagnitudeQuantity}
  \uses{def:physics:phyx-mini-0931:phyxminiproblems-problemphyxmini0931-arearatedimension}
  Magnitude of the instantaneous overlap-area change rate
\end{definition}
\begin{proof}
  This is the declared model definition of Area Rate Magnitude Quantity.
\end{proof}
\begin{definition}[Magnetic Flux Density Quantity]
  \label{def:physics:phyx-mini-0931:phyxminiproblems-problemphyxmini0931-magneticfluxdensityquantity}
  \lean{PhyXMiniProblems.ProblemPhyXMini0931.MagneticFluxDensityQuantity}
  \uses{def:physics:phyx-mini-0931:phyxminiproblems-problemphyxmini0931-magneticfluxdensitydimension}
  A nonnegative magnetic-flux-density magnitude
\end{definition}
\begin{proof}
  This is the declared model definition of Magnetic Flux Density Quantity.
\end{proof}
\begin{definition}[Magnetic Flux Magnitude Quantity]
  \label{def:physics:phyx-mini-0931:phyxminiproblems-problemphyxmini0931-magneticfluxmagnitudequantity}
  \lean{PhyXMiniProblems.ProblemPhyXMini0931.MagneticFluxMagnitudeQuantity}
  \uses{def:physics:phyx-mini-0931:phyxminiproblems-problemphyxmini0931-magneticfluxdimension}
  Magnitude of magnetic flux through the portion of loop in the field
\end{definition}
\begin{proof}
  This is the declared model definition of Magnetic Flux Magnitude Quantity.
\end{proof}
\begin{definition}[Magnetic Flux Rate Magnitude Quantity]
  \label{def:physics:phyx-mini-0931:phyxminiproblems-problemphyxmini0931-magneticfluxratemagnitudequantity}
  \lean{PhyXMiniProblems.ProblemPhyXMini0931.MagneticFluxRateMagnitudeQuantity}
  \uses{def:physics:phyx-mini-0931:phyxminiproblems-problemphyxmini0931-electromotiveforcedimension}
  Magnitude of a magnetic-flux change rate
\end{definition}
\begin{proof}
  This is the declared model definition of Magnetic Flux Rate Magnitude Quantity.
\end{proof}
\begin{definition}[Emf Magnitude Quantity]
  \label{def:physics:phyx-mini-0931:phyxminiproblems-problemphyxmini0931-emfmagnitudequantity}
  \lean{PhyXMiniProblems.ProblemPhyXMini0931.EmfMagnitudeQuantity}
  \uses{def:physics:phyx-mini-0931:phyxminiproblems-problemphyxmini0931-electromotiveforcedimension}
  Magnitude of an induced electromotive force
\end{definition}
\begin{proof}
  This is the declared model definition of Emf Magnitude Quantity.
\end{proof}
\begin{definition}[Resistance Quantity]
  \label{def:physics:phyx-mini-0931:phyxminiproblems-problemphyxmini0931-resistancequantity}
  \lean{PhyXMiniProblems.ProblemPhyXMini0931.ResistanceQuantity}
  \uses{def:physics:phyx-mini-0931:phyxminiproblems-problemphyxmini0931-electricalresistancedimension}
  A nonnegative electrical resistance
\end{definition}
\begin{proof}
  This is the declared model definition of Resistance Quantity.
\end{proof}
\begin{definition}[Current Magnitude Quantity]
  \label{def:physics:phyx-mini-0931:phyxminiproblems-problemphyxmini0931-currentmagnitudequantity}
  \lean{PhyXMiniProblems.ProblemPhyXMini0931.CurrentMagnitudeQuantity}
  \uses{def:physics:phyx-mini-0931:phyxminiproblems-problemphyxmini0931-electriccurrentdimension}
  Magnitude of an induced electric current
\end{definition}
\begin{proof}
  This is the declared model definition of Current Magnitude Quantity.
\end{proof}
\begin{definition}[nonnegative SIReadout]
  \label{def:physics:phyx-mini-0931:phyxminiproblems-problemphyxmini0931-nonnegativesireadout}
  \lean{PhyXMiniProblems.ProblemPhyXMini0931.nonnegativeSIReadout}
  Coherent-SI readout of a nonnegative dimensionful quantity
\end{definition}
\begin{proof}
  This is the declared model definition of nonnegative SIReadout.
\end{proof}
\begin{definition}[length In Meters]
  \label{def:physics:phyx-mini-0931:phyxminiproblems-problemphyxmini0931-lengthinmeters}
  \lean{PhyXMiniProblems.ProblemPhyXMini0931.lengthInMeters}
  \uses{def:physics:phyx-mini-0931:phyxminiproblems-problemphyxmini0931-lengthquantity, def:physics:phyx-mini-0931:phyxminiproblems-problemphyxmini0931-nonnegativesireadout}
  Metre readout of a physical length
\end{definition}
\begin{proof}
  This is the declared model definition of length In Meters.
\end{proof}
\begin{definition}[length In Centimeters]
  \label{def:physics:phyx-mini-0931:phyxminiproblems-problemphyxmini0931-lengthincentimeters}
  \lean{PhyXMiniProblems.ProblemPhyXMini0931.lengthInCentimeters}
  \uses{def:physics:phyx-mini-0931:phyxminiproblems-problemphyxmini0931-lengthquantity, def:physics:phyx-mini-0931:phyxminiproblems-problemphyxmini0931-lengthinmeters}
  Centimetre readout used by the two side-length guides
\end{definition}
\begin{proof}
  This is the declared model definition of length In Centimeters.
\end{proof}
\begin{definition}[time In Seconds]
  \label{def:physics:phyx-mini-0931:phyxminiproblems-problemphyxmini0931-timeinseconds}
  \lean{PhyXMiniProblems.ProblemPhyXMini0931.timeInSeconds}
  \uses{def:physics:phyx-mini-0931:phyxminiproblems-problemphyxmini0931-timequantity, def:physics:phyx-mini-0931:phyxminiproblems-problemphyxmini0931-nonnegativesireadout}
  Second readout of a physical clock time
\end{definition}
\begin{proof}
  This is the declared model definition of time In Seconds.
\end{proof}
\begin{definition}[speed In Meters Per Second]
  \label{def:physics:phyx-mini-0931:phyxminiproblems-problemphyxmini0931-speedinmeterspersecond}
  \lean{PhyXMiniProblems.ProblemPhyXMini0931.speedInMetersPerSecond}
  \uses{def:physics:phyx-mini-0931:phyxminiproblems-problemphyxmini0931-speedquantity}
  Metre-per-second readout of a speed
\end{definition}
\begin{proof}
  This is the declared model definition of speed In Meters Per Second.
\end{proof}
\begin{definition}[area In Square Meters]
  \label{def:physics:phyx-mini-0931:phyxminiproblems-problemphyxmini0931-areainsquaremeters}
  \lean{PhyXMiniProblems.ProblemPhyXMini0931.areaInSquareMeters}
  \uses{def:physics:phyx-mini-0931:phyxminiproblems-problemphyxmini0931-areaquantity, def:physics:phyx-mini-0931:phyxminiproblems-problemphyxmini0931-nonnegativesireadout}
  Square-metre readout of an area
\end{definition}
\begin{proof}
  This is the declared model definition of area In Square Meters.
\end{proof}
\begin{definition}[area Rate In Square Meters Per Second]
  \label{def:physics:phyx-mini-0931:phyxminiproblems-problemphyxmini0931-arearateinsquaremeterspersecond}
  \lean{PhyXMiniProblems.ProblemPhyXMini0931.areaRateInSquareMetersPerSecond}
  \uses{def:physics:phyx-mini-0931:phyxminiproblems-problemphyxmini0931-arearatemagnitudequantity, def:physics:phyx-mini-0931:phyxminiproblems-problemphyxmini0931-nonnegativesireadout}
  Square-metre-per-second readout of an area-rate magnitude
\end{definition}
\begin{proof}
  This is the declared model definition of area Rate In Square Meters Per Second.
\end{proof}
\begin{definition}[magnetic Flux Density In Teslas]
  \label{def:physics:phyx-mini-0931:phyxminiproblems-problemphyxmini0931-magneticfluxdensityinteslas}
  \lean{PhyXMiniProblems.ProblemPhyXMini0931.magneticFluxDensityInTeslas}
  \uses{def:physics:phyx-mini-0931:phyxminiproblems-problemphyxmini0931-magneticfluxdensityquantity, def:physics:phyx-mini-0931:phyxminiproblems-problemphyxmini0931-nonnegativesireadout}
  Tesla readout of magnetic flux density
\end{definition}
\begin{proof}
  This is the declared model definition of magnetic Flux Density In Teslas.
\end{proof}
\begin{definition}[magnetic Flux In Webers]
  \label{def:physics:phyx-mini-0931:phyxminiproblems-problemphyxmini0931-magneticfluxinwebers}
  \lean{PhyXMiniProblems.ProblemPhyXMini0931.magneticFluxInWebers}
  \uses{def:physics:phyx-mini-0931:phyxminiproblems-problemphyxmini0931-magneticfluxmagnitudequantity, def:physics:phyx-mini-0931:phyxminiproblems-problemphyxmini0931-nonnegativesireadout}
  Weber readout of a magnetic-flux magnitude
\end{definition}
\begin{proof}
  This is the declared model definition of magnetic Flux In Webers.
\end{proof}
\begin{definition}[magnetic Flux Rate In Webers Per Second]
  \label{def:physics:phyx-mini-0931:phyxminiproblems-problemphyxmini0931-magneticfluxrateinweberspersecond}
  \lean{PhyXMiniProblems.ProblemPhyXMini0931.magneticFluxRateInWebersPerSecond}
  \uses{def:physics:phyx-mini-0931:phyxminiproblems-problemphyxmini0931-magneticfluxratemagnitudequantity, def:physics:phyx-mini-0931:phyxminiproblems-problemphyxmini0931-nonnegativesireadout}
  Weber-per-second readout of a magnetic-flux-rate magnitude
\end{definition}
\begin{proof}
  This is the declared model definition of magnetic Flux Rate In Webers Per Second.
\end{proof}
\begin{definition}[emf Magnitude In Volts]
  \label{def:physics:phyx-mini-0931:phyxminiproblems-problemphyxmini0931-emfmagnitudeinvolts}
  \lean{PhyXMiniProblems.ProblemPhyXMini0931.emfMagnitudeInVolts}
  \uses{def:physics:phyx-mini-0931:phyxminiproblems-problemphyxmini0931-emfmagnitudequantity, def:physics:phyx-mini-0931:phyxminiproblems-problemphyxmini0931-nonnegativesireadout}
  Volt readout of an induced-emf magnitude
\end{definition}
\begin{proof}
  This is the declared model definition of emf Magnitude In Volts.
\end{proof}
\begin{definition}[resistance In Ohms]
  \label{def:physics:phyx-mini-0931:phyxminiproblems-problemphyxmini0931-resistanceinohms}
  \lean{PhyXMiniProblems.ProblemPhyXMini0931.resistanceInOhms}
  \uses{def:physics:phyx-mini-0931:phyxminiproblems-problemphyxmini0931-resistancequantity, def:physics:phyx-mini-0931:phyxminiproblems-problemphyxmini0931-nonnegativesireadout}
  Ohm readout of an electrical resistance
\end{definition}
\begin{proof}
  This is the declared model definition of resistance In Ohms.
\end{proof}
\begin{definition}[current Magnitude In Amperes]
  \label{def:physics:phyx-mini-0931:phyxminiproblems-problemphyxmini0931-currentmagnitudeinamperes}
  \lean{PhyXMiniProblems.ProblemPhyXMini0931.currentMagnitudeInAmperes}
  \uses{def:physics:phyx-mini-0931:phyxminiproblems-problemphyxmini0931-currentmagnitudequantity, def:physics:phyx-mini-0931:phyxminiproblems-problemphyxmini0931-nonnegativesireadout}
  Ampere readout of an electric-current magnitude
\end{definition}
\begin{proof}
  This is the declared model definition of current Magnitude In Amperes.
\end{proof}
\begin{definition}[Direction Vector]
  \label{def:physics:phyx-mini-0931:phyxminiproblems-problemphyxmini0931-directionvector}
  \lean{PhyXMiniProblems.ProblemPhyXMini0931.DirectionVector}
  A dimensionless direction vector in physical three-space
\end{definition}
\begin{proof}
  This is the declared model definition of Direction Vector.
\end{proof}
\begin{definition}[rightward Unit Vector]
  \label{def:physics:phyx-mini-0931:phyxminiproblems-problemphyxmini0931-rightwardunitvector}
  \lean{PhyXMiniProblems.ProblemPhyXMini0931.rightwardUnitVector}
  \uses{def:physics:phyx-mini-0931:phyxminiproblems-problemphyxmini0931-directionvector}
  Unit direction pointing right in the plane of the figure
\end{definition}
\begin{proof}
  This is the declared model definition of rightward Unit Vector.
\end{proof}
\begin{definition}[into Page Unit Vector]
  \label{def:physics:phyx-mini-0931:phyxminiproblems-problemphyxmini0931-intopageunitvector}
  \lean{PhyXMiniProblems.ProblemPhyXMini0931.intoPageUnitVector}
  \uses{def:physics:phyx-mini-0931:phyxminiproblems-problemphyxmini0931-directionvector}
  Unit direction pointing into the page
\end{definition}
\begin{proof}
  This is the declared model definition of into Page Unit Vector.
\end{proof}
\begin{definition}[horizontal Axis]
  \label{def:physics:phyx-mini-0931:phyxminiproblems-problemphyxmini0931-horizontalaxis}
  \lean{PhyXMiniProblems.ProblemPhyXMini0931.horizontalAxis}
  Horizontal coordinate used to describe the field half-plane
\end{definition}
\begin{proof}
  This is the declared model definition of horizontal Axis.
\end{proof}
\begin{definition}[Loop Vertex]
  \label{def:physics:phyx-mini-0931:phyxminiproblems-problemphyxmini0931-loopvertex}
  \lean{PhyXMiniProblems.ProblemPhyXMini0931.LoopVertex}
  The four vertices of the diamond-oriented square
\end{definition}
\begin{proof}
  This is the declared model definition of Loop Vertex.
\end{proof}
\begin{definition}[Loop Side]
  \label{def:physics:phyx-mini-0931:phyxminiproblems-problemphyxmini0931-loopside}
  \lean{PhyXMiniProblems.ProblemPhyXMini0931.LoopSide}
  The four equal sides of the square loop
\end{definition}
\begin{proof}
  This is the declared model definition of Loop Side.
\end{proof}
\begin{definition}[Guided Side]
  \label{def:physics:phyx-mini-0931:phyxminiproblems-problemphyxmini0931-guidedside}
  \lean{PhyXMiniProblems.ProblemPhyXMini0931.GuidedSide}
  The two sides carrying explicit 10 cm dimension guides in the raster
\end{definition}
\begin{proof}
  This is the declared model definition of Guided Side.
\end{proof}
\begin{definition}[Diagram Direction]
  \label{def:physics:phyx-mini-0931:phyxminiproblems-problemphyxmini0931-diagramdirection}
  \lean{PhyXMiniProblems.ProblemPhyXMini0931.DiagramDirection}
  Directions used for arrows and magnetic-field glyphs in the raster
\end{definition}
\begin{proof}
  This is the declared model definition of Diagram Direction.
\end{proof}
\begin{definition}[Diagram Color]
  \label{def:physics:phyx-mini-0931:phyxminiproblems-problemphyxmini0931-diagramcolor}
  \lean{PhyXMiniProblems.ProblemPhyXMini0931.DiagramColor}
  Stroke colors visible in the supplied raster
\end{definition}
\begin{proof}
  This is the declared model definition of Diagram Color.
\end{proof}
\begin{definition}[Page Normal Glyph]
  \label{def:physics:phyx-mini-0931:phyxminiproblems-problemphyxmini0931-pagenormalglyph}
  \lean{PhyXMiniProblems.ProblemPhyXMini0931.PageNormalGlyph}
  The cross/dot convention used for page-normal vectors
\end{definition}
\begin{proof}
  This is the declared model definition of Page Normal Glyph.
\end{proof}
\begin{definition}[Loop Shape]
  \label{def:physics:phyx-mini-0931:phyxminiproblems-problemphyxmini0931-loopshape}
  \lean{PhyXMiniProblems.ProblemPhyXMini0931.LoopShape}
  Geometric role assigned to the conducting loop
\end{definition}
\begin{proof}
  This is the declared model definition of Loop Shape.
\end{proof}
\begin{definition}[Square Orientation]
  \label{def:physics:phyx-mini-0931:phyxminiproblems-problemphyxmini0931-squareorientation}
  \lean{PhyXMiniProblems.ProblemPhyXMini0931.SquareOrientation}
  In-plane orientation of the square in the supplied image
\end{definition}
\begin{proof}
  This is the declared model definition of Square Orientation.
\end{proof}
\begin{definition}[Conductor Topology]
  \label{def:physics:phyx-mini-0931:phyxminiproblems-problemphyxmini0931-conductortopology}
  \lean{PhyXMiniProblems.ProblemPhyXMini0931.ConductorTopology}
  Topological role assigned to the conductor
\end{definition}
\begin{proof}
  This is the declared model definition of Conductor Topology.
\end{proof}
\begin{definition}[Translation Model]
  \label{def:physics:phyx-mini-0931:phyxminiproblems-problemphyxmini0931-translationmodel}
  \lean{PhyXMiniProblems.ProblemPhyXMini0931.TranslationModel}
  Motion model asserted by the problem statement
\end{definition}
\begin{proof}
  This is the declared model definition of Translation Model.
\end{proof}
\begin{definition}[Moving Square Loop Figure]
  \label{def:physics:phyx-mini-0931:phyxminiproblems-problemphyxmini0931-movingsquareloopfigure}
  \lean{PhyXMiniProblems.ProblemPhyXMini0931.MovingSquareLoopFigure}
  \uses{def:physics:phyx-mini-0931:phyxminiproblems-problemphyxmini0931-loopvertex, def:physics:phyx-mini-0931:phyxminiproblems-problemphyxmini0931-loopside, def:physics:phyx-mini-0931:phyxminiproblems-problemphyxmini0931-guidedside, def:physics:phyx-mini-0931:phyxminiproblems-problemphyxmini0931-diagramdirection, def:physics:phyx-mini-0931:phyxminiproblems-problemphyxmini0931-diagramcolor, def:physics:phyx-mini-0931:phyxminiproblems-problemphyxmini0931-pagenormalglyph}
  Literal presentation data transcribed from image 931.png. The displayed scalars are calibrated to independent physical quantities below; no induced current or answer-choice value is a field of this structure
\end{definition}
\begin{proof}
  This is the declared model definition of Moving Square Loop Figure.
\end{proof}
\begin{definition}[Moving Square Loop Setup]
  \label{def:physics:phyx-mini-0931:phyxminiproblems-problemphyxmini0931-movingsquareloopsetup}
  \lean{PhyXMiniProblems.ProblemPhyXMini0931.MovingSquareLoopSetup}
  \uses{def:physics:phyx-mini-0931:phyxminiproblems-problemphyxmini0931-lengthquantity, def:physics:phyx-mini-0931:phyxminiproblems-problemphyxmini0931-timequantity, def:physics:phyx-mini-0931:phyxminiproblems-problemphyxmini0931-speedquantity, def:physics:phyx-mini-0931:phyxminiproblems-problemphyxmini0931-areaquantity, def:physics:phyx-mini-0931:phyxminiproblems-problemphyxmini0931-arearatemagnitudequantity, def:physics:phyx-mini-0931:phyxminiproblems-problemphyxmini0931-magneticfluxdensityquantity, def:physics:phyx-mini-0931:phyxminiproblems-problemphyxmini0931-magneticfluxmagnitudequantity, def:physics:phyx-mini-0931:phyxminiproblems-problemphyxmini0931-magneticfluxratemagnitudequantity, def:physics:phyx-mini-0931:phyxminiproblems-problemphyxmini0931-emfmagnitudequantity, def:physics:phyx-mini-0931:phyxminiproblems-problemphyxmini0931-resistancequantity, def:physics:phyx-mini-0931:phyxminiproblems-problemphyxmini0931-currentmagnitudequantity, def:physics:phyx-mini-0931:phyxminiproblems-problemphyxmini0931-directionvector, def:physics:phyx-mini-0931:phyxminiproblems-problemphyxmini0931-loopshape, def:physics:phyx-mini-0931:phyxminiproblems-problemphyxmini0931-squareorientation, def:physics:phyx-mini-0931:phyxminiproblems-problemphyxmini0931-conductortopology, def:physics:phyx-mini-0931:phyxminiproblems-problemphyxmini0931-translationmodel, def:physics:phyx-mini-0931:phyxminiproblems-problemphyxmini0931-movingsquareloopfigure}
  The loop, field, geometry, clock, and independent electromagnetic observables. The maximum current is deliberately absent: inducedCurrentMagnitudeAt is an arbitrary time-dependent physical observable constrained only by later laws
\end{definition}
\begin{proof}
  This is the declared model definition of Moving Square Loop Setup.
\end{proof}
\begin{definition}[Matches Moving Square Loop Scenario]
  \label{def:physics:phyx-mini-0931:phyxminiproblems-problemphyxmini0931-matchesmovingsquareloopscenario}
  \lean{PhyXMiniProblems.ProblemPhyXMini0931.MatchesMovingSquareLoopScenario}
  \uses{def:physics:phyx-mini-0931:phyxminiproblems-problemphyxmini0931-rightwardunitvector, def:physics:phyx-mini-0931:phyxminiproblems-problemphyxmini0931-movingsquareloopsetup}
  Qualitative apparatus and motion roles stated by the problem
\end{definition}
\begin{proof}
  This is the declared model definition of Matches Moving Square Loop Scenario.
\end{proof}
\begin{definition}[Matches Supplied Moving Square Loop Figure]
  \label{def:physics:phyx-mini-0931:phyxminiproblems-problemphyxmini0931-matchessuppliedmovingsquareloopfigure}
  \lean{PhyXMiniProblems.ProblemPhyXMini0931.MatchesSuppliedMovingSquareLoopFigure}
  \uses{def:physics:phyx-mini-0931:phyxminiproblems-problemphyxmini0931-lengthincentimeters, def:physics:phyx-mini-0931:phyxminiproblems-problemphyxmini0931-speedinmeterspersecond, def:physics:phyx-mini-0931:phyxminiproblems-problemphyxmini0931-magneticfluxdensityinteslas, def:physics:phyx-mini-0931:phyxminiproblems-problemphyxmini0931-movingsquareloopsetup}
  Literal visual features and numerical labels from the primary raster
\end{definition}
\begin{proof}
  This is the declared model definition of Matches Supplied Moving Square Loop Figure.
\end{proof}
\begin{definition}[Matches Problem Measurements]
  \label{def:physics:phyx-mini-0931:phyxminiproblems-problemphyxmini0931-matchesproblemmeasurements}
  \lean{PhyXMiniProblems.ProblemPhyXMini0931.MatchesProblemMeasurements}
  \uses{def:physics:phyx-mini-0931:phyxminiproblems-problemphyxmini0931-lengthincentimeters, def:physics:phyx-mini-0931:phyxminiproblems-problemphyxmini0931-timeinseconds, def:physics:phyx-mini-0931:phyxminiproblems-problemphyxmini0931-speedinmeterspersecond, def:physics:phyx-mini-0931:phyxminiproblems-problemphyxmini0931-magneticfluxdensityinteslas, def:physics:phyx-mini-0931:phyxminiproblems-problemphyxmini0931-resistanceinohms, def:physics:phyx-mini-0931:phyxminiproblems-problemphyxmini0931-movingsquareloopsetup}
  Numerical data stated in the prose, including resistance and entry time
\end{definition}
\begin{proof}
  This is the declared model definition of Matches Problem Measurements.
\end{proof}
\begin{definition}[Has Physical Moving Loop Parameters]
  \label{def:physics:phyx-mini-0931:phyxminiproblems-problemphyxmini0931-hasphysicalmovingloopparameters}
  \lean{PhyXMiniProblems.ProblemPhyXMini0931.HasPhysicalMovingLoopParameters}
  \uses{def:physics:phyx-mini-0931:phyxminiproblems-problemphyxmini0931-lengthinmeters, def:physics:phyx-mini-0931:phyxminiproblems-problemphyxmini0931-speedinmeterspersecond, def:physics:phyx-mini-0931:phyxminiproblems-problemphyxmini0931-magneticfluxdensityinteslas, def:physics:phyx-mini-0931:phyxminiproblems-problemphyxmini0931-resistanceinohms, def:physics:phyx-mini-0931:phyxminiproblems-problemphyxmini0931-movingsquareloopsetup}
  Positivity and nondegeneracy conditions for the physical branch
\end{definition}
\begin{proof}
  This is the declared model definition of Has Physical Moving Loop Parameters.
\end{proof}
\begin{definition}[Models Uniform Into Page Half Plane Field]
  \label{def:physics:phyx-mini-0931:phyxminiproblems-problemphyxmini0931-modelsuniformintopagehalfplanefield}
  \lean{PhyXMiniProblems.ProblemPhyXMini0931.ModelsUniformIntoPageHalfPlaneField}
  \uses{def:physics:phyx-mini-0931:phyxminiproblems-problemphyxmini0931-magneticfluxdensityinteslas, def:physics:phyx-mini-0931:phyxminiproblems-problemphyxmini0931-intopageunitvector, def:physics:phyx-mini-0931:phyxminiproblems-problemphyxmini0931-horizontalaxis, def:physics:phyx-mini-0931:phyxminiproblems-problemphyxmini0931-movingsquareloopsetup}
  The supplied magnetic field occupies the half-plane to the right of the dashed boundary, is uniform there, vanishes outside, and points into the page. This uses Physlib's full space- and time-dependent magnetic-field object
\end{definition}
\begin{proof}
  This is the declared model definition of Models Uniform Into Page Half Plane Field.
\end{proof}
\begin{definition}[Satisfies Constant Translation Kinematics]
  \label{def:physics:phyx-mini-0931:phyxminiproblems-problemphyxmini0931-satisfiesconstanttranslationkinematics}
  \lean{PhyXMiniProblems.ProblemPhyXMini0931.SatisfiesConstantTranslationKinematics}
  \uses{def:physics:phyx-mini-0931:phyxminiproblems-problemphyxmini0931-lengthinmeters, def:physics:phyx-mini-0931:phyxminiproblems-problemphyxmini0931-timeinseconds, def:physics:phyx-mini-0931:phyxminiproblems-problemphyxmini0931-speedinmeterspersecond, def:physics:phyx-mini-0931:phyxminiproblems-problemphyxmini0931-movingsquareloopsetup}
  Constant-speed translation measured from the t = 0 entry event
\end{definition}
\begin{proof}
  This is the declared model definition of Satisfies Constant Translation Kinematics.
\end{proof}
\begin{definition}[Satisfies Diamond Boundary Sweep Geometry]
  \label{def:physics:phyx-mini-0931:phyxminiproblems-problemphyxmini0931-satisfiesdiamondboundarysweepgeometry}
  \lean{PhyXMiniProblems.ProblemPhyXMini0931.SatisfiesDiamondBoundarySweepGeometry}
  \uses{def:physics:phyx-mini-0931:phyxminiproblems-problemphyxmini0931-lengthinmeters, def:physics:phyx-mini-0931:phyxminiproblems-problemphyxmini0931-speedinmeterspersecond, def:physics:phyx-mini-0931:phyxminiproblems-problemphyxmini0931-areainsquaremeters, def:physics:phyx-mini-0931:phyxminiproblems-problemphyxmini0931-arearateinsquaremeterspersecond, def:physics:phyx-mini-0931:phyxminiproblems-problemphyxmini0931-movingsquareloopsetup}
  Geometry of a rigid square entering a vertical half-plane boundary. The vertical diagonal satisfies Pythagoras; every instantaneous swept width is at most this diagonal, and the bound is attained when the boundary passes through the top and bottom vertices. This law mentions area rate but neither emf nor current
\end{definition}
\begin{proof}
  This is the declared model definition of Satisfies Diamond Boundary Sweep Geometry.
\end{proof}
\begin{definition}[Satisfies Uniform Overlap Flux Law]
  \label{def:physics:phyx-mini-0931:phyxminiproblems-problemphyxmini0931-satisfiesuniformoverlapfluxlaw}
  \lean{PhyXMiniProblems.ProblemPhyXMini0931.SatisfiesUniformOverlapFluxLaw}
  \uses{def:physics:phyx-mini-0931:phyxminiproblems-problemphyxmini0931-areainsquaremeters, def:physics:phyx-mini-0931:phyxminiproblems-problemphyxmini0931-arearateinsquaremeterspersecond, def:physics:phyx-mini-0931:phyxminiproblems-problemphyxmini0931-magneticfluxdensityinteslas, def:physics:phyx-mini-0931:phyxminiproblems-problemphyxmini0931-magneticfluxinwebers, def:physics:phyx-mini-0931:phyxminiproblems-problemphyxmini0931-magneticfluxrateinweberspersecond, def:physics:phyx-mini-0931:phyxminiproblems-problemphyxmini0931-movingsquareloopsetup}
  For the uniform perpendicular field, flux magnitude is B times overlap area and flux-rate magnitude is B times swept-area-rate magnitude
\end{definition}
\begin{proof}
  This is the declared model definition of Satisfies Uniform Overlap Flux Law.
\end{proof}
\begin{definition}[Satisfies Faraday Induction Law]
  \label{def:physics:phyx-mini-0931:phyxminiproblems-problemphyxmini0931-satisfiesfaradayinductionlaw}
  \lean{PhyXMiniProblems.ProblemPhyXMini0931.SatisfiesFaradayInductionLaw}
  \uses{def:physics:phyx-mini-0931:phyxminiproblems-problemphyxmini0931-magneticfluxrateinweberspersecond, def:physics:phyx-mini-0931:phyxminiproblems-problemphyxmini0931-emfmagnitudeinvolts, def:physics:phyx-mini-0931:phyxminiproblems-problemphyxmini0931-movingsquareloopsetup}
  Magnitude form of Faraday's induction law, |ℰ| = |dΦ/dt|
\end{definition}
\begin{proof}
  This is the declared model definition of Satisfies Faraday Induction Law.
\end{proof}
\begin{definition}[Satisfies Ohms Law]
  \label{def:physics:phyx-mini-0931:phyxminiproblems-problemphyxmini0931-satisfiesohmslaw}
  \lean{PhyXMiniProblems.ProblemPhyXMini0931.SatisfiesOhmsLaw}
  \uses{def:physics:phyx-mini-0931:phyxminiproblems-problemphyxmini0931-emfmagnitudeinvolts, def:physics:phyx-mini-0931:phyxminiproblems-problemphyxmini0931-resistanceinohms, def:physics:phyx-mini-0931:phyxminiproblems-problemphyxmini0931-currentmagnitudeinamperes, def:physics:phyx-mini-0931:phyxminiproblems-problemphyxmini0931-movingsquareloopsetup}
  Ohm's law for the passive resistive loop, written without division as I R = |ℰ|. It is a general circuit law and contains no maximum value
\end{definition}
\begin{proof}
  This is the declared model definition of Satisfies Ohms Law.
\end{proof}
\begin{definition}[Is Maximum Readout]
  \label{def:physics:phyx-mini-0931:phyxminiproblems-problemphyxmini0931-ismaximumreadout}
  \lean{PhyXMiniProblems.ProblemPhyXMini0931.IsMaximumReadout}
  \uses{def:physics:phyx-mini-0931:phyxminiproblems-problemphyxmini0931-timequantity}
  A real readout is attained by, and bounds, a time-dependent observable
\end{definition}
\begin{proof}
  This is the declared model definition of Is Maximum Readout.
\end{proof}
\begin{definition}[Is Maximum Current In Amperes]
  \label{def:physics:phyx-mini-0931:phyxminiproblems-problemphyxmini0931-ismaximumcurrentinamperes}
  \lean{PhyXMiniProblems.ProblemPhyXMini0931.IsMaximumCurrentInAmperes}
  \uses{def:physics:phyx-mini-0931:phyxminiproblems-problemphyxmini0931-currentmagnitudeinamperes, def:physics:phyx-mini-0931:phyxminiproblems-problemphyxmini0931-movingsquareloopsetup, def:physics:phyx-mini-0931:phyxminiproblems-problemphyxmini0931-ismaximumreadout}
  A scalar ampere value is the maximum induced-current readout
\end{definition}
\begin{proof}
  This is the declared model definition of Is Maximum Current In Amperes.
\end{proof}
\begin{definition}[Answer Choice]
  \label{def:physics:phyx-mini-0931:phyxminiproblems-problemphyxmini0931-answerchoice}
  \lean{PhyXMiniProblems.ProblemPhyXMini0931.AnswerChoice}
  Labels attached to the four displayed integer-current choices
\end{definition}
\begin{proof}
  This is the declared model definition of Answer Choice.
\end{proof}
\begin{definition}[displayed Current In Amperes]
  \label{def:physics:phyx-mini-0931:phyxminiproblems-problemphyxmini0931-answerchoice-displayedcurrentinamperes}
  \lean{PhyXMiniProblems.ProblemPhyXMini0931.AnswerChoice.displayedCurrentInAmperes}
  \uses{def:physics:phyx-mini-0931:phyxminiproblems-problemphyxmini0931-answerchoice}
  Ampere value printed beside each answer label. The source strings are truncated after the numerals; amperes are inferred from the requested current
\end{definition}
\begin{proof}
  This is the declared model definition of displayed Current In Amperes.
\end{proof}
\begin{definition}[recorded Dataset Answer]
  \label{def:physics:phyx-mini-0931:phyxminiproblems-problemphyxmini0931-recordeddatasetanswer}
  \lean{PhyXMiniProblems.ProblemPhyXMini0931.recordedDatasetAnswer}
  \uses{def:physics:phyx-mini-0931:phyxminiproblems-problemphyxmini0931-answerchoice}
  The answer label recorded by the supplied dataset metadata
\end{definition}
\begin{proof}
  This is the declared model definition of recorded Dataset Answer.
\end{proof}
\begin{definition}[Is Closest Displayed Current]
  \label{def:physics:phyx-mini-0931:phyxminiproblems-problemphyxmini0931-isclosestdisplayedcurrent}
  \lean{PhyXMiniProblems.ProblemPhyXMini0931.IsClosestDisplayedCurrent}
  \uses{def:physics:phyx-mini-0931:phyxminiproblems-problemphyxmini0931-answerchoice}
  A displayed current is at least as close as every other displayed choice
\end{definition}
\begin{proof}
  This is the declared model definition of Is Closest Displayed Current.
\end{proof}
\begin{lemma}[vertical diagonal length formula]
  \label{lem:physics:phyx-mini-0931:phyxminiproblems-problemphyxmini0931-vertical-diagonal-length-formula}
  \lean{PhyXMiniProblems.ProblemPhyXMini0931.vertical_diagonal_length_formula}
  \uses{def:physics:phyx-mini-0931:phyxminiproblems-problemphyxmini0931-lengthinmeters, def:physics:phyx-mini-0931:phyxminiproblems-problemphyxmini0931-movingsquareloopsetup, def:physics:phyx-mini-0931:phyxminiproblems-problemphyxmini0931-hasphysicalmovingloopparameters, def:physics:phyx-mini-0931:phyxminiproblems-problemphyxmini0931-satisfiesdiamondboundarysweepgeometry}
  The vertical diagonal of a positive square is its side length times √2. This is a derived geometric result, not a field of the geometry premise
\end{lemma}
\begin{proof}
  The conclusion follows from the governing relations and figure readouts named in the statement of vertical diagonal length formula.
\end{proof}
\begin{lemma}[maximum overlap area growth rate]
  \label{lem:physics:phyx-mini-0931:phyxminiproblems-problemphyxmini0931-maximum-overlap-area-growth-rate}
  \lean{PhyXMiniProblems.ProblemPhyXMini0931.maximum_overlap_area_growth_rate}
  \uses{def:physics:phyx-mini-0931:phyxminiproblems-problemphyxmini0931-lengthinmeters, def:physics:phyx-mini-0931:phyxminiproblems-problemphyxmini0931-speedinmeterspersecond, def:physics:phyx-mini-0931:phyxminiproblems-problemphyxmini0931-arearateinsquaremeterspersecond, def:physics:phyx-mini-0931:phyxminiproblems-problemphyxmini0931-movingsquareloopsetup, def:physics:phyx-mini-0931:phyxminiproblems-problemphyxmini0931-hasphysicalmovingloopparameters, def:physics:phyx-mini-0931:phyxminiproblems-problemphyxmini0931-satisfiesdiamondboundarysweepgeometry, def:physics:phyx-mini-0931:phyxminiproblems-problemphyxmini0931-ismaximumreadout}
  The boundary sweep reaches maximum area rate v s √2
\end{lemma}
\begin{proof}
  The conclusion follows from the governing relations and figure readouts named in the statement of maximum overlap area growth rate.
\end{proof}
\begin{lemma}[maximum induced emf formula]
  \label{lem:physics:phyx-mini-0931:phyxminiproblems-problemphyxmini0931-maximum-induced-emf-formula}
  \lean{PhyXMiniProblems.ProblemPhyXMini0931.maximum_induced_emf_formula}
  \uses{def:physics:phyx-mini-0931:phyxminiproblems-problemphyxmini0931-lengthinmeters, def:physics:phyx-mini-0931:phyxminiproblems-problemphyxmini0931-speedinmeterspersecond, def:physics:phyx-mini-0931:phyxminiproblems-problemphyxmini0931-magneticfluxdensityinteslas, def:physics:phyx-mini-0931:phyxminiproblems-problemphyxmini0931-emfmagnitudeinvolts, def:physics:phyx-mini-0931:phyxminiproblems-problemphyxmini0931-movingsquareloopsetup, def:physics:phyx-mini-0931:phyxminiproblems-problemphyxmini0931-hasphysicalmovingloopparameters, def:physics:phyx-mini-0931:phyxminiproblems-problemphyxmini0931-satisfiesdiamondboundarysweepgeometry, def:physics:phyx-mini-0931:phyxminiproblems-problemphyxmini0931-satisfiesuniformoverlapfluxlaw, def:physics:phyx-mini-0931:phyxminiproblems-problemphyxmini0931-satisfiesfaradayinductionlaw, def:physics:phyx-mini-0931:phyxminiproblems-problemphyxmini0931-ismaximumreadout}
  Faraday's law turns the maximal swept-area rate into maximal emf
\end{lemma}
\begin{proof}
  The conclusion follows from the governing relations and figure readouts named in the statement of maximum induced emf formula.
\end{proof}
\begin{lemma}[maximum induced current formula]
  \label{lem:physics:phyx-mini-0931:phyxminiproblems-problemphyxmini0931-maximum-induced-current-formula}
  \lean{PhyXMiniProblems.ProblemPhyXMini0931.maximum_induced_current_formula}
  \uses{def:physics:phyx-mini-0931:phyxminiproblems-problemphyxmini0931-lengthinmeters, def:physics:phyx-mini-0931:phyxminiproblems-problemphyxmini0931-speedinmeterspersecond, def:physics:phyx-mini-0931:phyxminiproblems-problemphyxmini0931-magneticfluxdensityinteslas, def:physics:phyx-mini-0931:phyxminiproblems-problemphyxmini0931-resistanceinohms, def:physics:phyx-mini-0931:phyxminiproblems-problemphyxmini0931-movingsquareloopsetup, def:physics:phyx-mini-0931:phyxminiproblems-problemphyxmini0931-hasphysicalmovingloopparameters, def:physics:phyx-mini-0931:phyxminiproblems-problemphyxmini0931-satisfiesdiamondboundarysweepgeometry, def:physics:phyx-mini-0931:phyxminiproblems-problemphyxmini0931-satisfiesuniformoverlapfluxlaw, def:physics:phyx-mini-0931:phyxminiproblems-problemphyxmini0931-satisfiesfaradayinductionlaw, def:physics:phyx-mini-0931:phyxminiproblems-problemphyxmini0931-satisfiesohmslaw, def:physics:phyx-mini-0931:phyxminiproblems-problemphyxmini0931-ismaximumcurrentinamperes}
  Ohm's law gives the symbolic maximum-current formula
\end{lemma}
\begin{proof}
  The conclusion follows from the governing relations and figure readouts named in the statement of maximum induced current formula.
\end{proof}
% --- Archon declaration topology end ---
% --- Archon physics formalization source end ---
